\documentclass[11pt]{article}

\usepackage[margin=1in]{geometry}
\usepackage{graphicx}
\usepackage{booktabs}
\usepackage{array}
\usepackage{caption}
\usepackage{hyperref}
\usepackage{xcolor}
\usepackage{parskip}
\usepackage{textcomp}

\graphicspath{{../figures/}}

\hypersetup{
  colorlinks=true,
  urlcolor=blue,
  linkcolor=black
}

\title{Does Twitter feel the war in real time?\\
       \large VADER sentiment vs. Russia--Ukraine war events at CHTC scale}
\author{Nils Matteson \and Pravin Schmidley \and Will Tappa}
\date{STAT 405 / DSCP, April 2026}

\begin{document}
\maketitle

\section*{Introduction}

We ask whether public tweet sentiment shifts when a major
Russia--Ukraine war event happens. We scored \textbf{70.9 million tweets}
($\sim$52\,GB) from the Kaggle
\texttt{bwandowando/ukraine-russian-crisis-twitter-dataset} on CHTC
with HTCondor, then ran three analyses against 21 curated events:
Welch \textit{t}-tests, PELT change-point detection, and OLS
regression on event dummies with controls. The model explains
$\sim$9\% of daily variance ($R^2 = 0.09$, $F = 2.4$,
$p = 0.0007$); strongest effects are \textbf{Kherson liberated}
($\beta = +0.12$, $p = 0.006$) and \textbf{Kakhovka dam}
($\beta = -0.11$, $p = 0.02$). PELT flags 21 breaks; only one
matches a labeled event.

\section*{Data, computation, results}

\textbf{Data.} 480 daily CSVs ($\sim$52\,GB), Feb 2022 to June 2023.
Columns: \texttt{text}, \texttt{tweetcreatedts}, \texttt{language},
\texttt{retweetcount}, \texttt{followers}. Rows with empty text or
bad timestamps dropped ($<$2\%). After scoring: 70{,}876{,}101
tweets over 476 days. Primary pool is English (VADER is
English-trained).

\textbf{Computation.} An Apptainer container (pandas +
vaderSentiment, $\sim$340\,MB) was staged to
\texttt{/staging/nomatteson/}. Because \texttt{/staging} caps
users at 100 inodes, we packed the CSVs into 39 multi-day chunks.
For each chunk, \texttt{score\_tweets.py} streams the file in
200k-row blocks, runs VADER per tweet, and emits one row per
(date, language) with sums of compound, compound$^2$,
pos/neg/neu, followers, and followers$\times$compound --- sufficient
statistics, so the aggregator never re-reads raw tweets.
\textbf{Per job:} 1 CPU, $\sim$1.2\,GB RAM, 8\,GB disk; chunks
took 7--50\,min. \textbf{Total:} 39 parallel jobs, $\sim$1h40min
wall-clock. Raw tweets never left \texttt{/staging}.

\textbf{Results.}

\begin{figure}[h]
\centering
\includegraphics[width=\linewidth]{01_sentiment_timeline.png}
\caption{Daily VADER sentiment with annotated events and detected
change-points.}
\end{figure}

English daily mean compound sentiment drifts between $-0.05$ and
$-0.15$, with shorter negative dips around Bucha (April 2022) and
the deepest excursion during the Bakhmut period (spring 2023).
Of the 16 events with full $[-7, +7]$ windows, three pass
$|t| > 2$: Kharkiv counteroffensive ($+0.04$, $t = 2.7$),
annexation declared ($+0.02$, $t = 2.4$), and Bucha revealed
($-0.08$, $t = -2.2$).

\begin{figure}[h]
\centering
\includegraphics[width=\linewidth]{02_event_study.png}
\caption{Event study: post $-$ pre daily mean, 7-day windows.}
\end{figure}

The OLS regression
(\texttt{mean\_compound \textasciitilde\ t + log(n\_tweets) + 17 event dummies},
$n = 476$) gives $R^2 = 0.09$, $F = 2.4$ ($p = 0.0007$). Four
coefficients clear $|t| > 2$: Kherson liberated ($+0.12$,
$p = 0.006$), Kakhovka dam ($-0.11$, $p = 0.02$), anniversary
($+0.09$, $p = 0.04$), Zelensky at Congress ($+0.09$,
$p = 0.04$). PELT flags 21 breaks; only 2022-12-21 (Zelensky)
aligns with a labeled event. Reach-weighting shifts amplitude
but not sign.

\textbf{Weaknesses.} (i)~VADER is English-only; non-English
tweets look falsely neutral. (ii)~Twitter users are not the
public. (iii)~Durbin--Watson $= 0.70$ means autocorrelation
understates OLS standard errors. (iv)~Events cluster, so dummies
aren't causal. (v)~No bot filter.

\section*{Conclusion}

War events explain $\sim$9\% of daily variance on English
Twitter, but a handful register: Kherson's liberation and the
Kakhovka dam cross $|t| > 2$ with expected sign, and Bucha is
the clearest single-event dip. Most change-points don't match
labeled events --- either unlabeled drivers or VADER tracking
vocabulary drift. Future work: multilingual sentiment (XLM-R),
a bot filter, and regression-discontinuity around each event.

\subsection*{Contributions}

\begin{center}
\begin{tabular}{lcccc}
\toprule
Member & Proposal & Coding & Presentation & Report \\
\midrule
Nils Matteson    & 1 & 1 & 1 & 1 \\
Pravin Schmidley & 1 & 1 & 1 & 1 \\
Will Tappa       & 1 & 1 & 1 & 1 \\
\bottomrule
\end{tabular}
\end{center}

\noindent\textit{1 = full, 0.1--0.9 = partial, 0 = none.}

\vspace{0.5em}
\noindent Clone:
\texttt{git clone https://github.com/matteso1/STAT405project.git}

\end{document}
